\section{서론}\label{sec:intro}

여러 플레이어가 차례로 행동을 고르는 순차적 상호작용은 게임 트리로 나타낼 수 있다. 이것이 전개형 게임(extensive-form game, EFG)이다.
전개형 게임은 우연에 의한 사건과, 플레이어가 서로 구별하지 못하는 상태(불완전정보)를 함께 표현한다.
마르코프 결정 과정(MDP)이나 부분 관측 MDP(POMDP)와 구조가 비슷하지만, 여러 에이전트가 서로의 보상에 영향을 준다는 점이 다르다.
그래서 한 에이전트의 보상을 최대화하는 것이 목표가 아니라, 어느 플레이어도 혼자 전략을 바꿔 이득을 볼 수 없는 내쉬 균형을 찾는 것이 목표다.
유한 게임에는 혼합 전략 내쉬 균형이 항상 존재한다~\citep{nash1951non}.
2인 제로섬 게임에서는 최소최대 정리에 의해 균형 값이 하나로 정해지고, 균형 전략끼리 바꿔 써도 균형이 유지된다. 그래서 균형 하나를 찾는 것이 의미 있는 목표가 된다.
또 완전 기억(perfect recall) 게임에서는 행동 전략만으로 모든 혼합 전략을 표현할 수 있다~\citep{kuhn1953extensive}.

2인 제로섬 불완전정보 게임의 균형을 근사하는 표준 방법은 반사실적 후회 최소화(counterfactual regret minimization, CFR)다~\citep{zinkevich2007regret}.
CFR은 정보집합마다 "자신이 그 정보집합에 가려고 했다면"의 가치, 즉 상대와 우연의 도달확률만으로 가중한 반사실적 가치로 후회를 잰다.
이렇게 정의하면 전체 후회가 정보집합별 반사실적 후회의 합으로 위에서 묶인다. 그래서 정보집합마다 따로 후회를 줄이는 것만으로 전체 후회가 줄어든다.
정보집합마다 후회 매칭(regret matching)을 돌리면, 반복 전략들의 평균이 $O(1/\sqrt T)$의 exploitability로 균형에 다가간다.

그러나 CFR은 매 반복 게임 트리 전체를 순회한다. 트리는 깊이에 대해 지수적으로 커지므로, 큰 게임에서는 한 번의 순회도 어렵다.
MCCFR~\citep{lanctot2009monte}은 트리의 일부만 표본으로 순회하면서 반사실적 가치의 불편 추정량으로 후회를 갱신한다.
반복당 계산량은 훨씬 작으면서, 높은 확률로 CFR과 같은 차수의 수렴률을 가진다.
하지만 정보집합마다 따로 표를 두는 한, 처음 방문한 정보집합에는 다른 정보집합에서 배운 것을 줄 수 없다. 또 드물게 방문되는 정보집합의 전략은 표본 분산 때문에 크게 흔들린다.

그래서 큰 게임에서는 비슷한 정보집합이나 행동을 하나로 묶은 추상 게임을 푼다. 추상화에는 두 방향이 있다.
\begin{itemize}
  \item 연속적이거나 많은 행동을 대표 행동으로 줄이는 행동 추상화
  \item 비슷한 정보집합을 하나의 bucket으로 묶는 정보 추상화
\end{itemize}
둘 다 처음 보는 정보집합이나 행동에도 학습한 전략을 쓸 수 있게 한다. 이 논문은 정보 추상화를 다룬다.

고정된 정보 추상화에는 세 가지 문제가 있다.
\begin{enumerate}
  \item 추상 게임의 균형은 원 게임의 균형이 아니다. 고정 추상화에서 CFR이 수렴하는 곳은 추상 게임의 균형이고, 원 게임에서는 추상화 오차만큼 공략당한다.
  \item 추상화를 더 세밀하게 해도 원 게임에서 반드시 좋아지지는 않는다(추상화 병리 현상~\citep{waugh2009abstraction}).
    과거 정보를 잊는 불완전 기억 추상화에서는 CFR의 보장 자체가 일반적으로 성립하지 않는다~\citep{lanctot2012no}.
  \item 좋은 추상화를 만들려면 손패 세기의 군집처럼 게임마다 사람이 고른 특징이 필요하다~\citep{johanson2013evaluating,ganzfried2014potential}. 이런 특징은 일반 전개형 게임으로 옮기기 어렵다.
\end{enumerate}

CFR+IRA~\citep{cermak2020automated}는 학습 중에 거친 초기 추상화를 세분해 원 게임 균형으로 수렴시킨다.
하지만 매 반복 원 게임 트리 전체를 순회해야 해서, 순회 한 번도 어려운 큰 게임에는 쓸 수 없다.
저자들도 이것을 주된 한계로 꼽고, 표본 기법과의 결합을 후속 과제로 남겼다~\citep[6.1절]{cermak2020automated}.

\paragraph{이 논문의 방법.}
이 논문은 표본만으로 추상화를 세분하는 \emph{ASWISS}(Abstraction Split WIth Sampling by Strategy)를 제안한다. 본문은 가장 단순한 판인 \emph{순수 ASWISS}를 다룬다.
\begin{itemize}
  \item 게임 구조만으로 만든 가장 거친 분할(같은 플레이어, 같은 행동집합, 같은 시점)에서 시작한다.
  \item 현재 분할 위에서 외부 샘플링 MCCFR로 bucket 전략을 학습한다.
  \item 길이가 두 배씩 느는 epoch이 끝날 때마다, 각 정보집합에 그 epoch의 반복 전략들에 대한 \emph{누적} 반사실적 후회가 가장 큰 행동을 라벨로 붙인다. 그런 행동이 없으면 "없음" 라벨을 붙인다.
  \item 라벨이 갈린 bucket을 라벨별로 나눈다. 병합은 하지 않는다.
  \item 정보집합의 bucket은 정보집합별 표 없이, 시작 분할에서 분할 사건을 따라 내려가며 매번 다시 정한다(계보 라우팅). 그래서 처음 보는 정보집합에도 bucket이 정해진다.
\end{itemize}
이름의 "by Strategy"는 bucket을 정하는 것이 정보집합의 특징이 아니라 전략이 그 정보집합에 주는 업데이트 방향이라는 뜻이다.

\paragraph{기여.}
\begin{enumerate}
  \item \textbf{라벨 항등식과 수렴 정리.} bucket 안에서 상쇄되어 bucket이 느끼지 못하는 후회(숨은 후회)가 "모든 멤버에게 한 라벨을 강요할 때의 최소 비용"과 정확히 같음을 보인다(보조정리~\ref{lem:label}).
    이것으로 외부 샘플링 추정치만으로 분할을 판정해도, 분할이 유한 번이고 마지막 분할 이후 원 게임 NashConv가 0으로 감을 증명한다(정리~\ref{thm:epoch}).
    epoch의 반복 전략을 $K_k$개만 저장해 라벨을 재생하는 판에도 같은 꼴의 보장을 준다(정리~\ref{thm:replay}).
  \item \textbf{추상 게임의 수렴을 기다리는 변형과 정지 인증.} 추상 게임의 후회가 $\varepsilon$ 이하일 때만 판정하는 변형을 다룬다. 분할 없이 판정을 통과한 시점의 NashConv 인증(정리~\ref{thm:certificate})을 준다.
  \item \textbf{실험.} 정확한 exploitability를 잴 수 있는 작은 게임(Kuhn, Leduc, mini-stud)에서 CFR+IRA, 추상화 없는 CFR 계열, k-means 고정 추상화, 신경망 방법(MMD, ESCHER)과 비교한다.
    정보집합별 표 없이 계보를 따르는 비용도 실측한다.
  \item \textbf{파생 기법.} 계보 라우팅 비용과 방향 추정의 품질을 고치는 기법들을 같은 틀에서 정리한다(\ref{sec:derived}절). 평균 전략 방향과 실제 업데이트 방향의 정확한 차이(정리~\ref{thm:avgdir})가 그중 하나다.
\end{enumerate}

\paragraph{구성.}
2절은 관련 연구, 3절은 배경(CFR, MCCFR, CFR+IRA의 정리)이다.
4절에서 순수 ASWISS를 정의하고, 5절에서 이론 결과를 정리한다. 증명은 부록 A에 있다.
6절은 실험, 7절은 파생 기법, 8절은 한계와 후속 연구다.
